%% datacube-rs — Computers & Geosciences software paper
%% Elsevier elsarticle template. Compile: latexmk -pdf datacube-rs.tex
%%   (or: pdflatex; bibtex datacube-rs; pdflatex; pdflatex)
\documentclass[review]{elsarticle}

\usepackage{lineno,hyperref}
\usepackage{amsmath,amssymb}
\usepackage{booktabs}
\usepackage{graphicx}
\usepackage{xcolor}
\modulolinenumbers[5]

\journal{Computers \& Geosciences}

%% Numbered reference style (C&G)
\bibliographystyle{elsarticle-num}

\begin{document}

\begin{frontmatter}

\title{datacube-rs: An open-source Rust engine for streaming temporal
analysis of remote-sensing data cubes}

\author[usach]{Francisco Parra}
\ead{francisco.parra.o@usach.cl}
\address[usach]{Universidad de Santiago de Chile (USACH), Santiago, Chile}

\begin{abstract}
Open satellite archives have made remote sensing a temporal discipline, but the
tools that assemble image collections into space--time data cubes ship only
generic reducers, leaving trend, phenology and structural-break analysis to
separate single-purpose packages in different languages. We present
\texttt{datacube-rs}, an open-source engine written in Rust that unifies
temporal data-cube analysis with a validated set of per-pixel temporal
estimators in a single, portable artifact. It models a cube on the
$(\text{band}, y, x, \text{time})$ axes with the time dimension innermost,
enabling zero-copy per-pixel and per-chunk streaming parallelised with Rayon;
ingests satellite time series directly from STAC catalogues and Cloud-Optimized
GeoTIFF assets with cross-UTM-zone mosaicking; and provides ordinary-least-squares
trend, the Theil--Sen slope and tie-corrected Mann--Kendall test, harmonic
regression for seasonality, and OLS-CUSUM structural-break detection. Every
estimator is cross-validated to numerical parity against established reference
implementations (pyMannKendall, SciPy and statsmodels), passing 103 of 103
checks at a relative tolerance of $10^{-9}$. The same compiled core is delivered
as a native Rust library, a command-line tool, a Python extension (PyO3) and a
WebAssembly module that runs in the browser. We report parallel-scaling
benchmarks and a multiannual Sentinel-2 case study.
\end{abstract}

\begin{keyword}
Data cube \sep Satellite image time series \sep Rust \sep STAC \sep
Mann--Kendall \sep Phenology \sep Structural breaks
\end{keyword}

\end{frontmatter}

\linenumbers

\section{Introduction}
\label{sec:intro}

Open satellite archives have turned remote sensing into a fundamentally temporal
discipline. The full Landsat record and the Sentinel missions are now
distributed as free, analysis-ready data (ARD) through cloud object storage, and
a single study area may be covered by hundreds of observations per year. The
scientific questions that follow are, correspondingly, questions about time: is a
pixel's vegetation index trending up or down, and is that trend significant; when
does the growing season start and how is its amplitude changing; and where and
when does an abrupt disturbance---a fire, a clear-cut, a landslide---break the
expected seasonal signal. Answering them at scale requires both an efficient way
to assemble irregularly sampled, cloud-masked image collections into a regular
space--time cube, and a set of per-pixel statistical estimators applied along the
temporal axis.

A mature ecosystem addresses the first need. The \emph{gdalcubes} library
\citep{appel2019gdalcubes} builds on-demand regular data cubes from image
collections through GDAL, and endorses Cloud-Optimized GeoTIFF (COG) for
efficient partial reads. The Open Data Cube \citep{lewis2017agdc} provides a
Python framework for indexing and managing petabyte-scale ARD archives,
demonstrated on national high-performance infrastructure. The \emph{openEO} API
\citep{schramm2021openeo} harmonises heterogeneous cloud back-ends behind a
single, STAC-aligned virtual data-cube interface. \emph{xarray}
\citep{hoyer2017xarray} supplies the labeled $N$-dimensional array model that
underpins much of the scientific Python stack, and Google Earth Engine
\citep{gorelick2017gee} offers planetary-scale analysis to anyone with a browser.
Together these tools have made cube construction, ARD management and cloud access
broadly available.

The second need---temporal statistical analysis---is comparatively fragmented.
Cube engines expose generic reducers and aggregations (mean, median, percentile,
moving windows) but leave domain estimators to the user: gdalcubes, the Open Data
Cube, openEO and xarray ship no built-in monotonic-trend test, robust slope,
harmonic phenology model or structural-break detector, deferring these to
user-defined functions or external libraries. Google Earth Engine is a partial
exception---it provides Sen's slope and linear/robust regression reducers---but
it is a proprietary service whose computation cannot run offline and into whose
API existing workflows must be rewritten. Among cube tools, FORCE
\citep{frantz2019force} goes furthest, adding land-surface phenology and linear
trend analysis to a local ARD store, but it is a heavy C/C++ and
high-performance-computing stack that writes to a bespoke local grid rather than
streaming from cloud-native catalogues. Meanwhile the reference temporal methods
themselves live in separate, single-purpose packages and languages: the
Mann--Kendall trend-test family and Theil--Sen slope in \emph{pyMannKendall}
\citep{hussain2019pymannkendall}; the harmonic-plus-trend seasonality model in
Continuous Change Detection and Classification (CCDC; \citealp{zhu2014ccdc}); and
breakpoint detection on decomposed series in BFAST \citep{verbesselt2010bfast}. A
researcher who wants a significance-tested trend map, a phenology fit and a
disturbance-timing map from the same cube must therefore cross several tools,
languages and runtimes, and must trust that each per-pixel estimator is correctly
implemented.

Here we present \texttt{datacube-rs}, an open-source engine, written in Rust,
that unifies temporal data-cube analysis and a validated set of per-pixel
temporal estimators in a single, portable artifact. \texttt{datacube-rs} models a
cube on the $(\text{band}, y, x, \text{time})$ axes with the time dimension
innermost, so that each pixel's series is a contiguous slice and can be
streamed---per pixel or by spatial chunk---without copying; per-pixel computation
is parallelised across cores with Rayon. It ingests scenes directly from STAC
catalogues and COG assets (reusing the cloud stack of the SurtGIS engine),
mosaicking across UTM zones onto a common grid. On top of the cube it provides
ordinary-least-squares trend, the Theil--Sen robust slope and the tie-corrected
Mann--Kendall test, harmonic (Fourier) regression for seasonality and phenology,
and structural-break detection by an OLS-CUSUM test with recursive binary
segmentation---the latter inspired by BFAST's decompose-then-detect strategy,
though it deliberately uses a different test (OLS-CUSUM rather than BFAST's
OLS-MOSUM with Bai--Perron selection; see Section~\ref{sec:methods}). Every
estimator treats missing values as \texttt{NaN} and drops them pairwise, and
every reported quantity is cross-validated to numerical parity against
established reference implementations---pyMannKendall, SciPy and
statsmodels---passing 103 of 103 checks at a relative tolerance of
$10^{-9}$. The same compiled core is exposed as a native Rust library, a
command-line tool, a Python extension (via PyO3) and a WebAssembly module that
runs the estimators in the browser.

The contributions of this paper are:
\begin{itemize}
  \item a streaming $(\text{band}, y, x, \text{time})$ cube model whose
  contiguous time axis enables zero-copy per-pixel and per-chunk iteration,
  parallelised with Rayon;
  \item a set of per-pixel temporal estimators---OLS trend, Theil--Sen,
  tie-corrected Mann--Kendall, harmonic regression and OLS-CUSUM structural
  breaks---implemented natively and validated to $10^{-9}$ numerical parity
  against pyMannKendall, SciPy and statsmodels;
  \item direct ingestion of satellite time series from STAC/COG with
  cross-UTM-zone mosaicking, plus temporal compositing and gap-filling on the
  cube;
  \item a single dependency-light Rust core delivered as four deployment
  targets---library, CLI, Python and in-browser WebAssembly---with reproducible
  benchmarks of its parallel scaling.
\end{itemize}

The remainder of the paper is organised as follows.
Section~\ref{sec:methods} states the temporal estimators and their reference
implementations. Section~\ref{sec:engine} describes the \texttt{datacube-rs}
engine: the cube data model, the streaming iterator and chunking, STAC/COG
ingestion and the deployment targets. Section~\ref{sec:validation} reports the
numerical-parity validation, performance benchmarks and a real Sentinel-2 case
study. Section~\ref{sec:discussion} discusses how \texttt{datacube-rs} relates to
existing tools and its current limitations, and Section~\ref{sec:conclusions}
concludes.

\section{Temporal estimators}
\label{sec:methods}
% TODO: OLS trend; Theil-Sen + tie-corrected Mann-Kendall (Var(S) formula,
% cite Kendall 1975 / Mann 1945); harmonic regression (CCDC model, Eq. 3);
% OLS-CUSUM break test (Brownian-bridge p-value) + binary segmentation.
% State the OLS-MOSUM (BFAST) vs OLS-CUSUM distinction explicitly.

\section{The datacube-rs engine}
\label{sec:engine}
% TODO: Fig. 1 architecture. 3.1 cube model; 3.2 streaming iterator + chunking
% (Rayon); 3.3 STAC/COG ingest + cross-zone mosaic; 3.4 deployment targets.

\section{Validation and benchmarking}
\label{sec:validation}
% TODO: 4.1 numerical parity (103/103, table); 4.2 criterion benchmarks
% (scaling); 4.3 Sentinel-2 case study (Santiago / UTM 18-19 boundary).

\section{Discussion}
\label{sec:discussion}
% TODO: positioning vs gdalcubes/ODC/openEO/GEE/FORCE; what Rust single-binary
% + WASM buys; limitations + roadmap (GeoZarr backing, compositing).

\section{Conclusions}
\label{sec:conclusions}
% TODO

%% --- Back-matter (Elsevier / C&G order) ---
\section*{CRediT authorship contribution statement}
\textbf{Francisco Parra:} Conceptualization, Methodology, Software, Validation,
Writing -- original draft.

\section*{Code availability section}
\noindent
Name of the code/library: \texttt{datacube-rs}\\
Contact: Francisco Parra, francisco.parra.o@usach.cl\\
Hardware requirements: any x86-64/ARM64 machine; benchmarks on a 16-core laptop\\
Program language: Rust (with Python bindings via PyO3 and a WebAssembly target)\\
Software required: Rust toolchain ($\geq$1.94, edition 2024)\\
Program size: TODO\\
Source code: \url{https://github.com/franciscoparrao/datacube-rs}\\
License: MIT OR Apache-2.0\\
Archived version: TODO (Zenodo DOI)

\section*{Data availability}
The Sentinel-2 case study uses freely available data from the Microsoft
Planetary Computer STAC catalogue. The synthetic series and scripts used for the
numerical-parity validation are included in the source repository.

\section*{Declaration of competing interest}
The author declares no competing interests.

\section*{Acknowledgements}
% TODO

\bibliography{refs}

\end{document}
